% Responsable: Tahiny
\subsection{Sistema embebido de una planta}

\textbf{Alcance del caso.} Se asume que una planta industrial dedicada a la
fabricación y envasado de productos necesita
implementar un sistema embebido que permita
automatizar y supervisar el funcionamiento de su
línea de producción. Actualmente, algunos procesos
requieren intervención manual, lo que puede
ocasionar retrasos, errores operativos, desperdicio
de materiales y dificultades para detectar fallos
en la maquinaria.
El sistema estará basado en un controlador embebido industrial que
recibirá información de sensores de temperatura,
presión, proximidad y nivel, con el propósito de
monitorear las condiciones de operación en tiempo
real. A partir de estos datos, el controlador
accionará automáticamente equipos como motores,
bandas transportadoras, bombas y válvulas,
siguiendo una secuencia de funcionamiento
previamente establecida. Asimismo, contará con una interfaz de supervisión
que permitirá a los operarios visualizar el estado
de los equipos, consultar parámetros de
funcionamiento y recibir alertas ante situaciones
anormales. Se contempla el registro de eventos
y la ejecución de acciones de protección, como
detener de manera controlada determinados equipos
cuando se detecten fallos.

\textbf{Límites del alcance:}
El proyecto se limita al monitoreo, control y
automatización de la línea de producción, sin
incluir procesos administrativos, comerciales
o financieros de la planta.
Se asume que los requisitos principales de
funcionamiento y seguridad serán definidos antes
del desarrollo, debido a la criticidad de las
operaciones industriales.

\subsubsection{Caracterización}
\subsection{Cliente, actores y características del proyecto}
\begin{description}
  \item[Cliente y objetivo.] Se asume que el cliente es la administración
  de la planta y que busca supervisar la operación de una línea productiva,
  reducir fallos de operación y registrar estados de los equipos.

  \item[Actores.] Se asume la participación de una jefatura de producción
  (acepta objetivos operativos), operadores (usan las señales y alarmas),
  personal de mantenimiento (diagnostica fallos), especialistas en
  automatización (integran sensores y actuadores) y responsables de
  seguridad industrial (revisan peligros y medidas de protección).

  \item[Estabilidad de requisitos.] \textbf{Alta en el núcleo, bajo
  supuesto:} se presume que la función de control y las interfaces físicas
  pueden definirse antes de construir. Los ajustes de parámetros o la
  presentación de alarmas podrían cambiar y requieren gestión de cambios.

  \item[Criticidad.] \textbf{Alta, bajo supuesto:} una orden incorrecta a
  un actuador o una lectura defectuosa podría afectar la continuidad del
  proceso, provocar daños en equipos o generar peligros. Esta valoración
  exige revisar y probar los comportamientos ante fallos; no presupone
  que todo el software cumpla una función de seguridad certificada.

  \item[Contrato y regulación.] Se asume un contrato o acuerdo técnico con
  alcance, interfaces y criterios de aceptación definidos. La normativa
  obligatoria \textbf{no está indicada en el enunciado} y deberá
  determinarse según país, maquinaria y funciones implementadas. Si hay
  funciones eléctricas/electrónicas/programables relacionadas con la
  seguridad, corresponde evaluar la pertinencia de IEC 61508-1:2010,
  sin afirmar que sea automáticamente obligatoria
  \cite{iec61508_2010}.

  \item[Disponibilidad del cliente.] \textbf{Limitada, bajo supuesto:}
  se presume que operadores y técnicos trabajan por turnos. Serían
  posibles reuniones programadas para aclarar requisitos, revisar
  documentos y participar en pruebas de aceptación, pero no necesariamente
  reuniones diarias.

  \item[Tamaño e integración.] \textbf{Medio, bajo supuesto:} el alcance
  incluye firmware o lógica de control, sensores, actuadores, una interfaz
  de operación y pruebas de integración con el equipo físico. No se
  dispone de estimaciones reales de esfuerzo ni de componentes.
\end{description}

\subsubsection{Modelo elegido y justificación}
\textbf{Elección propuesta: modelo en cascada con verificaciones técnicas
tempranas y control de cambios.} La elección solo se mantiene si la
revisión inicial confirma los supuestos anteriores. Se justifica por:
\begin{enumerate}
  \item \textbf{Estabilidad del núcleo de control:} permitiría acordar
  previamente qué señales se leen, qué acciones se autorizan y cuáles son
  los límites de operación, disminuyendo cambios tardíos en interfaces.
  \item \textbf{Criticidad y trazabilidad:} permite establecer requisitos
  y criterios de prueba identificables, relacionarlos con el diseño y
  conservar actas y resultados para la aceptación.
  \item \textbf{Participación del cliente por hitos:} se adapta mejor a
  revisiones programadas con producción, mantenimiento y seguridad
  industrial que a la exigencia de decisiones operativas diarias.
\end{enumerate}

\textbf{Por qué no el iterativo como modelo principal:} sería preferible
si existieran dudas relevantes sobre sensores, interfaces o respuestas del
controlador que solo pudieran resolverse mediante prototipos. Bajo los
supuestos actuales, esa incertidumbre se considera manejable con ensayos
preliminares, sin redefinir todo el proceso por ciclos.

\textbf{Por qué no el ágil como modelo principal:} las posibles ventanas
reducidas de acceso a la planta y la necesidad de acordar previamente las
interfaces y evidencias de aceptación hacen menos conveniente depender de
retroalimentación continua. Esto no significa que las prácticas ágiles
sean incompatibles con sistemas embebidos.

\subsubsection{Riesgos del modelo elegido y mitigación}
% Ronda 3 (sábado, hasta 22:30): riesgos concretos y cómo mitigarlos.
